\chapter*{Introduction}
\addcontentsline{toc}{chapter}{Introduction}

Arduino is a very popular microcomputer among amateurs and professionals alike. Despite this, testing the code for an Arduino machine can be fairly difficult. Debugging options are limited and the user needs to rely on additional measuring equipment. The usual choice is to test the code with an AVR emulator. Emulation is the process of running machine code on a hardware that doesn't match its instruction set. AVR is the processor that controls the Arduino board. These emulators are sometimes part of a bigger logical circuit simulation system, other times they are just a single piece of software.

Goal of this work is to skip the emulation process and create a working Arduino simulator by mimicking the changes of digital values on the board's pinout. The user code will be compiled directly to the x86-64 architecture instruction set, which is our target platform, and thus able to run at the native speed of the host machine.

Bundled with this simulator is a graphic application for creating and testing logical circuits, with Arduino as it's main component. The point of this application is to attach to the running Arduino simulator and visualize the expected Arduino behavior with user's code on a given circuit. We want to give the user the power to compose their own circuits, and define the behavior and looks of their components. Additionally, the user will be able to debug their program, while the GUI is running, using their preferred IDE and C++ debugger, providing real-time feedback and inspection when stepping through the code.

\section{Arduino programming guide}

\subsection{Sketches}

The Arduino programs are commonly referred to as \textit{sketches}. A sketch is not a single file, but an entire directory, usually populated by files with the \texttt{.ino} extension, contaning the source code of the program. \texttt{.cpp}, \texttt{.c}, \texttt{.S} (assembly files), and header files like \texttt{.h} are also accepted as files contaning source code. Every sketch must contain a \texttt{.ino} file with a file name matching the sketch root directory name \cite{arduino_sketch_specs}.

The Arduino programming language, in which \texttt{.ino} files are written, is a C++ dialect with very subtle differences, mostly concerning linking of files and ommitting function declarations. Bodies of function definitions and variable declarations can be considered as a perfectly valid C/C++ code. 

The essential functions that an Arduino programmer has to define in their sketch are \texttt{void setup()} and \texttt{void loop()}. The \texttt{setup()} function is called once, when the program is booted. That is when the Arduino board is powered up, the program is uploaded to it, or when the reset button is pressed. The \texttt{loop()} function is called in a loop indefinitely. This is the proper place to define any interations with the board, detecting inputs, updating timers, etc.

\subsection{Arduino board interface}

The target board \texttt{Arduino Uno R3 SMD} provides 6 analog and 14 digital general purpose pins, some of which can be used for I2C and SPI protocol communication, or PWM modulation. The most simple to use and simulate are the digital pins for general purpose I/O, so this thesis focuses exlucively on using these pins as such \cite{arduino_r3_smd_specs}.

The digital pins can be manipulated using only these 3 functions:
\begin{itemize}
    \item \texttt{void pinMode(pin, mode)} - configures specified pin either as an input or an output
    \item \texttt{value digitalRead(pin)} - reads the voltage on the pin, stores it as a digitla value, either 1 or 0
    \item \texttt{void digitalWrite(pin, value)} - when set to output, sets the voltage of the pin to corresponding value. If set to input, current driven components may not work correctly, since the current is suppressed by an internal pull-up resistor. 
\end{itemize}

Pins are identified by their index. \texttt{mode} and \texttt{value}'s underlying type is also an integer, but it is preferred to use predefined constants, such as \texttt{INPUT}/\texttt{OUTPUT}, \texttt{HIGH}/\texttt{LOW} and such. Below is an example of a sketch from the Arduino's documentation, using only the described functions:

\begin{lstlisting}[language=C, caption=Sketch example \cite{arduino_digital_read}]
int ledPin = 13;  // LED connected to digital pin 13
int inPin = 7;    // pushbutton connected to digital pin 7
int val = LOW;    // variable to store the read value

void setup() {
  pinMode(ledPin, OUTPUT);  // sets the digital pin 13 as output
  pinMode(inPin, INPUT);    // sets the digital pin 7 as input
}

void loop() {
  val = digitalRead(inPin);   // read the input pin
  digitalWrite(ledPin, val);  // sets LED to the button's value
}
\end{lstlisting}

\subsection{Build process}

Before the code is compiled, all \texttt{.ino} go through a simple preprocessing stage. \texttt{\#line} directives are added. These link to the original line and file, when the code is debugged. All files are concatenated together, function declarations are generated at the top, and the \texttt{Arduino.h} is included. This header file is crucial as it declares the \texttt{void} and \texttt{loop} functions, functions for manipulating pins, common constants and mathematical functions, etc. At this point the sketch becomes a complete C++ source code.

It is also at this where we steer away from the standard compilation process, and apply our own \texttt{Arduino.h} implementation in our simulator. Usually, the C++ code would be linked to an implementation that contains binaries and constants for the given microchip installed on the target board. Then it would be compiled with a compiler for the given instruction set (e.g. \texttt{avr-gcc}). The code would be then transmitted over USB to the user's Arduino.

In our case, we include the entire preprocessed sketch in our simulator project written in C++, and compile it as whole for the instruction set of the host machine. This executable will implement mock functions declared in the \texttt{Arduino.h}, arduino's internal state, and a way to communicate the change of state to an external GUI application.

The entire process from preprocessing to compilation, can be orchestrated with the use of \texttt{arduino-cli} command line application. We will use it to preprocess the \texttt{.ino} for us, so that we end up with a clean C++ file we can include in our project. 

\section{Specific goals of the thesis}
